\subsection{Scope classification}
\label{sec:domain-classification}

Independently of the two-axis availability judgment above, every paper is also
checked for whether it is genuinely in scope for this study at all --- a separate
judgment from the arXiv category under which a paper was retrieved
(Section~\ref{sec:category-corpus-construction}), or from any keyword signal found
near it. Neither of those signals constitutes genuine content judgment: retrieval
reflects only which category a paper was filed under, and the full-text availability
scan (Section~\ref{sec:candidate-list-gen}) only checks for substring presence, not
whether a paper is actually about the topic in question.

This matters specifically on the quantum side, where \texttt{quant-ph} as an arXiv
category covers all of quantum physics, not just quantum computing ---
quantum sensing, quantum foundations, quantum optics, and quantum networking all live
there too, alongside quantum computing itself. A paper using quantum hardware to
detect a gravitational wave, for instance, is not a quantum computing paper merely
because it involves qubits: the qubit is acting as a sensor, not executing a
computation, and the paper's contribution is a physical measurement, not an advance in
computing capability (see Section~\ref{sec:results-detail} for further discussion of
this distinction). The category-based corpus-construction procedure
(Section~\ref{sec:category-corpus-construction}) is deliberately kept broad
(maximizing recall, per the same design principle already applied to the full-text
availability scan in Section~\ref{sec:candidate-list-gen}) rather than pre-filtered
for this distinction at pull time, so this check is deferred to a dedicated
classification step instead of narrowing the query and losing volume.

Scope classification is performed with the same reading depth and per-paper judgment
as the availability classification (Section~\ref{sec:framework}), not a reduced or
heuristic pass: an agent reads a paper's extracted text (or digest, at full-corpus
scale, per Section~\ref{sec:extraction-classification}) and judges whether the paper
is genuinely quantum computing research, or falls into one of the specific
non-computing categories already established earlier in this project for exactly this
problem: \emph{Atomic Physics}, \emph{Quantum Sensing}, \emph{Quantum Magnonics},
\emph{Quantum Optomechanics}, \emph{Photonic Quantum Computing}, \emph{Trapped-Ion
Computing}, \emph{Neutral-Atom Computing}, \emph{Semiconductor Spin Qubits}, and
\emph{High Energy Physics}, plus a general \emph{Not Quantum Computing} fallback for a
paper that does not fit any of the above. A paper assigned one of these labels is
retained in the corpus (so that corpus size and candidate-generation yield remain
auditable) but excluded from every subsequent quantum-computing-specific statistic in
this report. No equivalent out-of-scope category set is currently defined on the CS
side; whether CS's broad pull similarly needs an out-of-scope carve-out remains an
open question.

For papers not yet classified for availability, this judgment is obtained in the same
pass as the Axis A/Axis B classification --- the same read of the paper's text yields
both judgments, avoiding a second full read-through of the corpus. For the 3{,}472
quantum papers already classified for availability before this scope-classification
methodology was adopted (Table~\ref{tab:corpus-stats}), a separate backfill pass adds
only the scope judgment, at the same rigor as if it had been obtained in the combined
pass; their existing Axis A/Axis B classification, evidence, and confidence fields are
left unmodified. This backfill is implemented as its own reusable script, applying the
same deterministic-extraction-plus-agent-classification pattern already established
for availability classification, so that scope classification can be independently
re-run in the future without requiring the availability judgment to be redone. As of
this writing, this methodology has been designed, but neither the backfill pass over
the 3{,}472 already-classified quantum papers nor the combined-pass classification of
the CS corpus has yet been executed.

\subsection{Scope note}
% NOTE: this subsection previously contrasted the hand-reviewed pilot table against
% full-corpus progress; the pilot table has since been removed in favor of the
% aggregate-results blueprint (Section~\ref{sec:results-aggregate}), so most of the
% original point here is now moot. Left in only for the still-relevant provenance/
% cross-validation caveat below -- consider deleting this subsection entirely.
The extraction-verification-classification pipeline (Section~\ref{sec:extraction-classification})
is intended for scaling classification to the full corpus in each domain
(Table~\ref{tab:corpus-stats}). As of this writing, it has been run against a
3{,}478-paper quantum sample, with 3{,}472 of those papers classified
(Table~\ref{tab:corpus-stats}); this is, in either case, a single-pass review that has
not yet been cross-validated by a second reviewer.
